%%%%%%%%%%%%%%%%%%%%%part.tex%%%%%%%%%%%%%%%%%%%%%%%%%%%%%%%%%%
% 
% sample part title
%
% Use this file as a template for your own input.
%
%%%%%%%%%%%%%%%%%%%%%%%% Springer %%%%%%%%%%%%%%%%%%%%%%%%%%

\begin{partbacktext}
\part{Preliminaries Overview}
\noindent 

% Short introduction - need to expand but not more than one page

These first two chapters cover context for the rest of the book. Chapter \ref{chp:introduction} is an introduction to clinical physiology including what it means to take a muscle centered approach. It then reviews the core concepts of physiology. 

Chapter \ref{chp:fundamentals} includes the broader context for muscle fibers. Specifically, how they connect together and some of the characteristics they display when working together as muscles. Physical therapists are interested in muscle fibers because when connected together they form muscles which create movement. 
\end{partbacktext}